
\section{Introduction}\label{intro}

\subsection{Statement of  the results}  

About 15 years ago Yoccoz proved that the Julia set of at most finitely many renormalizable quadratic polynomial
$f: z\mapsto z^2+c$ with all periodic points repelling is locally connected (see \cite{H,M}).
In this paper, we generalize this result to higher degree unicritical polynomials: 

\proclaim Theorem A. 
The Julia set $J(f)$ of 
at most finitely renormalizable unicritical polynomial  $f: z\mapsto  z^d+c$  with all periodic points repelling 
is locally connected.

This result follows from {\it a priori} bounds in an appropriate 
``Modified Principle Nest''  of puzzle pieces,
$$
  E^0\Supset E^1\Supset\dots \ni 0: 
$$

\proclaim Theorem B. The (modified) principal moduli stay away from zero:
$$
    \mod(E^{i-1}\sm E^i)\geq \mu > 0.
$$

These {\it a priori} bounds imply that the puzzle pieces $E^i$ shrink to the critical point,
which yields Theorem A by a standard argument.%


\subsection{Techniques}

As usual in holomorphic dynamics, our proof has two sides: combinatorial and analytic.
Our combinatorial tool is a refined  Principal Nest techniques of \cite{puzzle}, 
while the analytic tool is a recently established
Quasi-Invariance Law (Covering Lemma) in conformal geometry \cite{covering lemma}.  
Let us briefly comment on both sides. 

The puzzle machinery was introduced to holomorphic dynamics by Branner and Hubbard \cite{BH}
(in the context of cubic polynomials with one escaping critical point) and
Yoccoz \cite{H,M} (in the context of quadratic polynomials). 
The idea is to tile shrinking neighborhoods
of the Julia set into topological disks called {\it puzzle pieces}, 
and to translate the dynamics on $J(f)$ to the combinatorics of these tilings.

An efficient way to describe these combinatorics is given by
the {\it Principal nest} of puzzle pieces around the origin,
$V^0\supset V^1\supset\dots V^n\dots\ni 0$, which is inductively constructed so that 
the first return maps $f^{n_i}: V^{i}\ra V^{i-1}$ are unicritical branched coverings
 \cite{puzzle}. 
It turns out that this nest is not quite suitable for our purposes, 
so we modify it slightly to obtain  
a {\it dynasty of kingdom map}, see \S \ref{principal nest}. 

We then observe that since the return times in the dynasty grow exponentially,
one can send some puzzle piece $E^{i-1}$ 
to the top level by an appropriate composition $\Psi$ of the kingdom maps, 
while the next puzzle piece, $E^i$, will go at most five levels up ({\it time inequality}).
Thus, the map $\Psi|\, E^i$ has a bounded degree, which 
puts us in a position to apply the analytic techniques of \cite{covering lemma}.
 
The puzzle bears  complete information about the Julia set only if
the puzzle pieces shrink to points, so this is a key geometric issue of the theory.
To handle this issue,  Branner \& Hubbard and Yoccoz made use of the 
Series Law from conformal geometry.%
\footnote{also called  the   Gr\"otzsch Inequality.}
It was immediately realized, however, that this method would not work for higher degree polynomials,
so that in the higher degree case the problem has remained open since then.

A new analytic tool that we exploit is a Covering Lemma (Quasi-Invariance Law) in conformal geometry \cite{covering lemma}
which roughly asserts that given a branched covering $g: U\ra V$ of degree $N$ which restricts to a branched covering
$g: A\ra B$ of degree $d$ such that $\mod (U\sm A) $ is small (depending on $N$), 
then, under a certain ``Collar Assumption'',  $ \mod(V\sm B)$  is  comparable with  $d^2 \mod(U\sm A)$  
(independently of $N$) -- see \S \ref{covering lemma sec} for the precise statement.    

The Covering Lemma allows us to transfer moduli information from deep levels of the dynasty to 
shallow ones, and to argue that if on some deep levels the moduli are small, then they must be even 
smaller on  shallow ones. This certainly implies that, in fact, the moduli can never  be too small (Theorem B).

\ssk
Note in conclusion that
for real $c$, Theorem A  was proved before by Levin \& van Strien \cite{LS}.
The method  used in \cite{LS} exploited real symmetry in a substantial way. 

\ssk In the forthcoming notes (joint with A. Avila and W. Shen)
      our {\it a priori} bounds will be used to prove {\it rigidity} of the 
      unicritical polynomials under consideration.

\subsection{Terminology and Notation}   

  A {\it topological disk} means a simply connected  domain in $\C$.
 
  We let $\orb(z)\equiv \orb_g(z)= (g^n z)_{n=0}^\infty$ be the {\it orbit} of $z$ under a map $g$.

  Given a map $g: U\ra V$ and a domain $D\subset V$, 
components of $g^{-1}(D)$ are called {\it pullbacks}  of $D$ under $g$. 
Given a connected set $X\subset g^{-1}(D)$, we let $g^{-1} (D)|X$ be the pullback of $D$
containing $X$.    

Given a subset $W\subset V$, the {\it first landing map} $H$ to $W$ is defined (on the set of points $z$
whose orbits intersect $W$) as follows:  $H(z)= f^l z$, where $l\geq 0$ is the first moment for which
$f^l z\in W$.

We say that a map $g: U\ra V$ is {\it unicritical} if it has one critical point
(of arbitrary local degree)


\msk {\bf Acknowledgment.}
 We thank Artur Avila for careful reading the manuscript and making a number of useful comments.
We also thank all the Foundations that have supported this work: 
the Guggenheim Fellowship, Clay Mathematics Institute, NSF, and NSERC.
